\title{DeepSeekMath: Pushing the Limits of Mathematical Reasoning in Open Language Models}

\begin{abstract}
The intricacies and inherent structure of mathematical reasoning present notable difficulties for language models. In this work, we present DeepSeekMath~7B, a model developed by extending the pre-training of DeepSeek-Coder-Base-v1.5 7B with an additional 120B math-centric tokens acquired from Common Crawl, complemented by both natural language and code corpora. Without the assistance of auxiliary toolkits or voting mechanisms, DeepSeekMath~7B attains a remarkable 51.7\% accuracy on the advanced MATH benchmark, narrowing the performance gap with state-of-the-art models such as Gemini-Ultra and GPT-4. Moreover, applying self-consistency sampling across 64 generations enables DeepSeekMath~7B to reach a 60.9\% score on MATH. This notable progress in mathematical reasoning is largely credited to two crucial strategies: (1) exploiting the vast potential embedded in open web resources via a carefully constructed data selection pipeline, and (2) proposing Group Relative Policy Optimization (GRPO), a modification of Proximal Policy Optimization (PPO) that not only augments mathematical reasoning capacity but also enhances PPO’s memory efficiency.
\end{abstract}

\section{Introduction}

Large language models (LLMs) have fundamentally transformed mathematical reasoning within artificial intelligence, leading to notable improvements on benchmarks for quantitative \citep{MATH} as well as geometry reasoning \citep{trinh2024solving}. These models have also demonstrated value in aiding humans with the resolution of challenging mathematical tasks \citep{tao}. Nonetheless, state-of-the-art systems like GPT-4 \citep{gpt4} and Gemini-Ultra \citep{gemini} remain closed to the public, and existing open-source alternatives exhibit substantially lower performance.

In this work, we present DeepSeekMath, a specialized language model that exhibits superior mathematical reasoning compared to open-source alternatives and nearly reaches the proficiency of GPT-4 across standard academic benchmarks. Our approach centers on the construction of the DeepSeekMath Corpus, a comprehensive and high-quality pre-training dataset containing 120B math tokens. This corpus is compiled from Common Crawl (CC) by leveraging a classifier built with fastText \citep{joulin2016fasttext}. Initially, this classifier is trained with positive samples from OpenWebMath \citep{openwebmath}, complemented by a wide range of unrelated web content as negative samples. Utilizing this model, we proceed to extract further positive data from the CC, which undergoes a meticulous human annotation process for additional refinement. The resulting curated dataset is used to further retrain the classifier, thereby enhancing its effectiveness. According to evaluation outcomes, the scale and quality of the corpus are validated, as the DeepSeekMath-Base 7B model attains a score of 64.2\% on GSM8K \citep{gsm8k} and 36.2\% on the MATH competition dataset \citep{MATH}, surpassing Minerva 540B \citep{minerva}. Moreover, the multilingual composition of the DeepSeekMath Corpus enables notable gains in performance on Chinese mathematical evaluations \citep{wei2023cmath,agieval}. We propose that our methods for mathematical data curation and processing offer valuable insights for the field, with ample opportunity for further advancements.

DeepSeekMath-Base builds upon DeepSeek-Coder-Base-v1.5 7B \citep{deepseek-coder}, given that initializing from a model trained on code yields superior outcomes to using a general-purpose LLM. In addition, our analysis shows that training on mathematical data also boosts performance on benchmarks such as MMLU \citep{mmlu} and BBH \citep{bbh}, suggesting gains not only in mathematical proficiency but also in broader reasoning skills.

Subsequent to pre-training, we conduct mathematical instruction tuning on DeepSeekMath-Base utilizing data featuring chain-of-thought \citep{cot}, program-of-thought \citep{pot,pal}, and reasoning incorporating external tools \citep{tora}. The final model, DeepSeekMath-Instruct 7B, outperforms all other 7B models and achieves performance levels comparable to 70B open-source instruction-tuned models.

In addition, we propose Group Relative Policy Optimization (GRPO), a novel variant of reinforcement learning (RL) derived from Proximal Policy Optimization (PPO) \citep{schulman2017proximal}. Unlike conventional PPO, GRPO dispenses with the critic network, estimating the baseline instead from group-wise scores, which leads to a significant reduction in training cost. Utilizing only a portion of English instruction-tuning data, GRPO achieves marked performance gains compared to the already competitive DeepSeekMath-Instruct—showing improvements both in in-domain tasks (GSM8K: 82.9\% $\rightarrow$ 88.2\%, MATH: 46.8\% $\rightarrow$ 51.7\%) and out-of-domain mathematical benchmarks (e.g., CMATH: 84.6\% $\rightarrow$ 88.8\%) within the reinforcement learning process. Furthermore, we establish a cohesive framework for interpreting various techniques, including Rejection Sampling Fine-Tuning (RFT) \citep{yuan2023scaling}, Direct Preference Optimization (DPO) \citep{dpo}, PPO, and our proposed GRPO. Through this unified perspective, we observe that these algorithms may be interpreted as variations or simplifications of RL strategies. Comprehensive experiments—spanning online versus offline training, outcome versus process-level supervision, as well as single-turn versus multi-turn RL—are conducted to rigorously examine key properties of this paradigm. Lastly, we articulate the reasons underlying the effectiveness of our RL-based approach in enhancing instruction-tuned models, and outline future avenues for advancing RL methods under this unified framework.

\subsection{Contributions}
Our contribution includes scalable math pre-training, along with the exploration and analysis of reinforcement learning.

\textbf{Math Pre-Training at Scale}

\begin{itemize}[topsep=0pt]
    \item This study demonstrates that the openly available Common Crawl dataset holds significant utility for mathematical tasks. Through the deployment of a carefully structured data selection process, we assemble the DeepSeekMath Corpus, a rigorously curated dataset comprising 120B tokens sourced from web pages emphasizing mathematical material. This corpus surpasses the size of the mathematics-related web content used by Minerva \citep{minerva} by nearly sevenfold, and exceeds the scale of the recently introduced OpenWebMath \citep{openwebmath} by a factor of nine.

\item The DeepSeekMath-Base 7B model, after pre-training, demonstrates similar levels of performance to Minerva 540B \citep{minerva}, suggesting that the parameter count alone does not solely determine mathematical reasoning proficiency. Notably, smaller models trained on curated, high-quality datasets are also capable of exhibiting robust performance.

\item We present results from our experiments involving mathematical training. Pre-training models on code enhances their performance on mathematical tasks, regardless of whether they use tools. These results help address the longstanding question: \textit{does code training improve reasoning abilities?} Our evidence suggests that, specifically for mathematical reasoning, it indeed does.

    \item Although training on arXiv papers is common, especially in many math-related papers, it brings no notable improvements on all mathematical benchmarks adopted in this paper.
\end{itemize}

\textbf{Exploration and Analysis of Reinforcement Learning}

\begin{itemize}[topsep=0pt]
    \item We present Group Relative Policy Optimization (GRPO), a reinforcement learning algorithm that is both efficient and effective. Unlike conventional approaches such as Proximal Policy Optimization (PPO), GRPO eliminates the need for a critic model by leveraging group score-based baseline estimates, thereby substantially lowering the computational cost during training.
    \item We provide empirical evidence that applying GRPO leads to substantial improvements in the instruction-tuned model DeepSeekMath-Instruct, utilizing only instruction-tuning datasets. In addition, we report gains in generalization to out-of-domain data throughout the reinforcement learning phase.
    \item We propose a cohesive framework that connects various methods, including RFT, DPO, PPO, and GRPO. Our extensive experimental analysis investigates key components of this unified paradigm, covering areas such as online versus offline training modalities, supervision based on outcomes versus processes, and comparisons between single-turn and iterative reinforcement learning strategies.
    \item Drawing on the proposed paradigm, we examine the underlying mechanisms contributing to the effectiveness of reinforcement learning and outline several promising directions for enhancing the reinforcement learning of LLMs in future work.
\end{itemize}

\subsection{Summary of Evaluations and Metrics}

\begin{itemize}[topsep=0pt]
    \item \textbf{English and Chinese Mathematical Reasoning}: We perform extensive evaluations of our models using both English and Chinese benchmark datasets that encompass mathematical questions ranging from primary to collegiate levels. For English-language benchmarks, we include GSM8K \citep{gsm8k}, MATH \citep{MATH}, SAT \citep{llemma}, OCW Courses \citep{minerva}, and MMLU-STEM \citep{mmlu}. On the Chinese side, the evaluation employs MGSM-zh \citep{mgsm}, CMATH \citep{wei2023cmath}, Gaokao-MathCloze \citep{agieval}, as well as Gaokao-MathQA \citep{agieval}. Our assessment covers both the production of fully articulated solutions in natural language without external tool assistance and the models' proficiency in resolving problems through Python-based computations.

In evaluations on English benchmarks, DeepSeekMath-Base matches the performance of the proprietary Minerva 540B model \citep{minerva}, and exceeds all available open-source base models—such as Mistral 7B \citep{mistral} and Llemma-34B \citep{llemma}—irrespective of whether these models have received math-specific pre-training, and frequently achieves this with a notable lead. Importantly, DeepSeekMath-Base also attains superior outcomes on Chinese benchmarks. This may be attributed to our decision not to limit math pre-training data exclusively to English, in contrast with previous approaches \citep{minerva,llemma}, and instead incorporate high-quality mathematical data from non-English sources. Further enhancement through mathematical instruction tuning and reinforcement learning yields DeepSeekMath-Instruct and DeepSeekMath-RL, both of which showcase robust results—most notably, surpassing the 50\% accuracy threshold for the first time on the competition-level MATH dataset among open-source models.

\item \textbf{Formal Mathematics}: We assess DeepSeekMath-Base on the informal-to-formal theorem proving benchmark detailed in \citep{dsp_proof}, employing miniF2F \citep{minif2f} with Isabelle \citep{isabelle} as the selected proof assistant. Our results indicate that DeepSeekMath-Base achieves notable performance in few-shot autoformalization tasks.

\item \textbf{Natural Language Understanding, Reasoning, and Code}: To thoroughly evaluate the model’s abilities in general comprehension, reasoning, and programming, we test DeepSeekMath-Base using the Massive Multitask Language Understanding (MMLU) suite \citep{mmlu}, which spans 57 diverse multiple-choice questions, as well as the BIG-Bench Hard (BBH) benchmark \citep{bbh}, comprised of 23 difficult tasks that typically necessitate multi-step reasoning. In addition, HumanEval \citep{codex} and MBPP \citep{mbpp} serve as standard assessments for coding proficiency. Pre-training on mathematical data enhances both linguistic and reasoning-related outcomes.

\section{Math Pre-Training}

\subsection{Data Collection and Decontamination} This section details the methodology for assembling the DeepSeekMath Corpus utilizing data from Common Crawl. Illustrated in Figure \ref{fig:data_collect}, our iterative pipeline facilitates the organized collection of an extensive mathematical text corpus by beginning with a seed corpus—namely, a compact yet high-quality math-focused dataset. Importantly, this strategy can be extended to encompass additional domains, including but not limited to coding.

We begin by selecting OpenWebMath \citep{openwebmath}—a curated collection of high-quality mathematics-related web documents—as our foundational seed corpus. From this dataset, we train a fastText classifier \citep{joulin2016fasttext} to identify and retrieve additional mathematical web pages resembling those in OpenWebMath. Our training set consists of 500,000 positive samples randomly drawn from the seed corpus and 500,000 negative samples taken from Common Crawl. Training is conducted using an open-source implementation\footnote{\url{https://fasttext.cc}}, with parameters set as follows: vector dimension of 256, learning rate of 0.1, maximum word n-gram length of 3, minimum word frequency threshold of 3, and three training epochs. Prior to extraction, we reduce the raw Common Crawl dataset via URL-based and near-duplicate deduplication, narrowing it down to 40B HTML documents. Utilizing the trained fastText model, we retrieve pages from this deduplicated corpus that are most similar to mathematical content. Subsequently, these candidates are ranked according to their fastText-derived scores, and only those attaining the highest rankings are retained, effectively filtering out lower quality content. To determine the appropriate amount of data to keep, we conduct pre-training runs on subsets consisting of the top 40B, 80B, 120B, and 160B tokens by score; for this initial phase, we opt to retain the top 40B tokens.

Following the initial round of data collection, a significant portion of mathematical web pages remains unacquired. This is primarily attributed to the limited variety present in the positive examples used to train the fastText model. To address this issue and enhance the model's effectiveness, we seek out additional mathematical web sources in order to diversify and expand the seed corpus. Our method involves segmenting the complete Common Crawl dataset into separate domains, where each domain is defined as all web pages that originate from a common base URL. For every domain, we compute the proportion of its web pages captured during the first collection phase. Domains with more than 10\% of their pages retrieved are categorized as being related to mathematics (for example, \url{mathoverflow.net}).

Next, we carry out manual annotation of the URLs within these domains that are indicative of mathematical content, such as \url{mathoverflow.net/questions}. Any web pages accessible through these annotated URLs, but which have not yet been gathered, are incorporated into the seed corpus. This iterative methodology allows us to broaden the pool of positive examples, leading to an enhanced fastText model that achieves higher recall of mathematical web pages in future iterations. After completing four rounds of data collection using this strategy, we assemble a total of 35.5M mathematical web pages, amounting to 120B tokens. Upon reaching the fourth round, we observe that approximately 98\% of the eventual dataset was already acquired in the third round, prompting us to halt the data collection process at this stage.

In order to prevent benchmark contamination, we adopt the approach described by \cite{deepseek-coder}, removing web pages that feature questions or answers from established English mathematical benchmarks, including GSM8K \citep{gsm8k} and MATH \citep{MATH}, as well as Chinese benchmarks such as CMATH \citep{wei2023cmath} and AGIEval \citep{agieval}. Our filtering method stipulates that any passage containing a 10-gram sequence exactly matching any sub-string from these evaluation benchmarks is excluded from the math training dataset. For benchmark passages that are shorter than 10 grams but contain at least 3 grams, we apply an exact matching strategy to detect and eliminate any web pages suspected of contamination.

\subsection{Validating the Quality of the DeepSeekMath~Corpus}

We run pre-training experiments to investigate how the DeepSeekMath~Corpus is compared with the recently released math-training corpora:

\begin{itemize}[topsep=0pt]
    \item \textbf{MathPile} \citep{mathpile}: This dataset comprises 8.9B tokens collected from diverse resources such as textbooks, Wikipedia, ProofWiki, CommonCrawl, StackExchange, and arXiv, with arXiv constituting over 85\% of the total data;
    \item \textbf{OpenWebMath} \citep{openwebmath}: A 13.6B token dataset built by extracting mathematical material from CommonCrawl through content filtering;
    \item \textbf{Proof-Pile-2} \citep{llemma}: Encompassing OpenWebMath, AlgebraicStack (10.3B tokens of mathematical programming code), and arXiv articles (28.0B tokens), this mathematical corpus adopts the same arXiv:Web:Code proportion of 2:4:1 as reported in \cite{llemma} during experiments.
\end{itemize}

\subsubsection{Training Setting}
\label{sec:quality-policy}
Our approach involves fine-tuning a base language model with 1.3B parameters, which follows the architectural design of the DeepSeek LLMs \citep{deepseek-llm}, henceforth referred to as DeepSeek-LLM 1.3B, using mathematical datasets. Each distinct mathematical dataset is used to train an individual model, with each undergoing training for 150B tokens. The entirety of the training process employs the resource-efficient HAI-LLM framework \citep{haillm}. Consistent with DeepSeek LLMs’ established methodology, we utilize the AdamW optimizer \citep{adamW} with parameters $\beta_1=0.9$, $\beta_2=0.95$, and a $\mathrm{weight\_decay}$ of $0.1$. The learning rate schedule is structured such that the rate ramps up to its maximum after 2,000 warmup steps, reduces to 31.6\% of this maximum after 80\% of total steps, and drops further to 10.0\% of the maximum by the time 90\% of training steps have elapsed. The peak learning rate is set at $5.3\times10^{-4}$. We configure training with a batch size totaling 4M tokens and use a context window of 4K tokens.

\subsubsection{Evaluation Results}

\textbf{The DeepSeekMath~Corpus is of high quality, covers multilingual mathematical content, and is the largest in size.}

\begin{itemize}[topsep=0pt]
    \item \textbf{High-quality}:
    We assess the model’s downstream abilities using few-shot chain-of-thought prompting \cite{cot} across eight mathematical benchmarks. The results in Table \ref{tab:corpora_comparison} reveal that models trained with the DeepSeekMath~Corpus consistently outperform those using other datasets. Furthermore, as illustrated in Figure \ref{fig:corpora_comparisons}, after processing 50B tokens—equivalent to a complete epoch over Proof-Pile-2—the model trained on the DeepSeekMath Corpus surpasses its Proof-Pile-2 counterpart, highlighting the superior average quality of the DeepSeekMath Corpus. 
    \item \textbf{Multilingual}:
    Incorporating materials in multiple languages, the DeepSeekMath~Corpus primarily features English and Chinese content. Table \ref{tab:corpora_comparison} demonstrates that using the DeepSeekMath~Corpus for training yields substantial gains in mathematical reasoning in both English and Chinese. In comparison, existing mathematical corpora, which largely focus on English, contribute limited benefit and may adversely affect Chinese language mathematical capabilities.
    \item \textbf{Large-scale}:
    Compared to previous mathematical corpora, the DeepSeekMath~Corpus is significantly larger in size. Figure \ref{fig:corpora_comparisons} depicts the DeepSeek-LLM 1.3B model’s steeper and more sustained learning improvements when trained on DeepSeekMath~Corpus. On the other hand, the smaller baseline corpora have already undergone several training passes, with their corresponding models quickly approaching a performance ceiling.
\end{itemize}

\subsection{Training and Evaluating DeepSeekMath-Base 7B}

This section presents DeepSeekMath-Base 7B, a foundational model demonstrating advanced reasoning capabilities, particularly in mathematical contexts. The model's parameters are initialized from DeepSeek-Coder-Base-v1.5 7B \citep{deepseek-coder} and it undergoes training on a corpus containing 500B tokens. The data sources are apportioned as follows: 56\% originates from the DeepSeekMath Corpus, 4\% from AlgebraicStack, 10\% from arXiv, 20\% is code obtained from Github, while the remaining 10\% comprises natural language content drawn from Common Crawl in both English and Chinese. The training configuration generally follows the protocol described in Section \ref{sec:quality-policy}, with two key exceptions: the maximum learning rate is set to 4.2e-4 and a batch size of 10M tokens is utilized.

In this study, we systematically evaluate the mathematical proficiency of DeepSeekMath-Base 7B by examining its performance in generating complete mathematical solutions independently, utilizing tools to address mathematical problems, and engaging in rigorous formal theorem proving. Additionally, we extend our analysis to capture a broader overview of the base model's capabilities, encompassing its aptitude for natural language comprehension, logical reasoning, and programming abilities.

\paragraph{Mathematical Problem Solving with Step-by-Step Reasoning}

We assess the mathematical problem-solving capabilities of DeepSeekMath-Base by employing few-shot chain-of-thought prompting \citep{cot} on a suite of eight benchmarks presented in both English and Chinese. These evaluation datasets span a range of quantitative reasoning tasks—such as GSM8K \citep{gsm8k}, MATH \citep{MATH}, and CMATH \citep{wei2023cmath}—as well as multiple-choice formats exemplified by MMLU-STEM \citep{mmlu} and Gaokao-MathQA \citep{agieval}. Collectively, these benchmarks address mathematical content from basic to advanced collegiate levels, and represent diverse mathematical domains.

As indicated in Table \ref{tab:base_math_cot}, DeepSeekMath-Base 7B achieves the highest scores among all open-source base models on each of the eight evaluated benchmarks. This group includes notable models such as the popular Mistral 7B \citep{mistral} and the recently introduced Llemma 34B \citep{llemma}, the latter of which was specifically trained on the Proof-Pile-2 dataset \citep{llemma} for mathematics. Of particular significance, DeepSeekMath-Base outperforms all open-source base models on the challenging MATH competition dataset by more than 10 percentage points. Furthermore, it exceeds the performance of Minerva 540B \citep{minerva}—a much larger, closed-source base model built on top of PaLM \citep{palm} and additionally tuned with mathematical corpora—despite being only a fraction of its size.

\paragraph{Mathematical Problem Solving with Tool Use} We assess mathematical reasoning aided by program execution on the GSM8K and MATH benchmarks by employing few-shot program-of-thought prompting, following the method described by \citet{pot,pal}. For each problem, the model generates a Python script—which may use libraries such as \textit{math} and \textit{sympy}—to carry out complex calculations. The answer is taken from the program's output upon execution. Table \ref{tab:base_math_pot_proof} demonstrates that DeepSeekMath-Base 7B surpasses the previous leading model, Llemma 34B.

\paragraph{Formal Mathematics}

The automation of formal proof plays a significant role in promoting the precision and dependability of mathematical reasoning, as well as streamlining the proving process, a topic that has seen growing scholarly interest recently. In this study, we assess the performance of DeepSeekMath-Base 7B on the informal-to-formal proving task introduced in \citep{dsp_proof}. This task involves constructing a formal proof given an informal description, its formal representation, and an informal justification. Our evaluation utilizes miniF2F \citep{minif2f}, which serves as a testbed for formalization in Olympiad-level mathematics; for each question, we employ few-shot prompting to produce corresponding Isabelle formal proofs. Consistent with the methodology outlined in \cite{dsp_proof}, we use large language models to compose initial proof outlines and then rely on the readily available automated prover Sledgehammer \citep{sledgehammer} to address any omitted arguments. Table \ref{tab:base_math_pot_proof} highlights that DeepSeekMath-Base 7B achieves notable results in automating formal proof construction.

\paragraph{Natural Language Understanding, Reasoning, and Code}

We assess the natural language understanding capabilities of models using the MMLU benchmark \citep{mmlu}, reasoning abilities using BBH \citep{bbh}, and coding proficiency with HumanEval \citep{codex} and MBPP \citep{mbpp}. According to Table \ref{tab:understanding_reasoning_code}, DeepSeekMath-Base 7B demonstrates notable improvements on both MMLU and BBH compared to its predecessor, DeepSeek-Coder-Base-v1.5 \citep{deepseek-coder}, which underscores the beneficial effects of mathematical training on understanding and reasoning in language tasks. Moreover, by incorporating code tokens during continued training, DeepSeekMath-Base 7B preserves the coding benchmark performances previously achieved by DeepSeek-Coder-Base-v1.5. On the whole, DeepSeekMath-Base 7B achieves considerably better results than the general-purpose Mistral 7B \citep{mistral} across the evaluated reasoning and coding benchmarks.

\section{Supervised Fine-Tuning}

\subsection{SFT Data Curation}

A comprehensive mathematical instruction-tuning dataset is developed, comprising problems in both English and Chinese across multiple areas of mathematics and spanning a range of difficulty levels. Each problem is annotated with a solution employing chain-of-thought (CoT) \citep{cot}, program-of-thought (PoT) \citep{pot,pal}, or tool-integrated reasoning methodologies \citep{tora}. The resulting training set contains a total of 776K examples.

\begin{itemize}[topsep=0pt]
    \item \textbf{English mathematical datasets}: We provide tool-integrated solution annotations for GSM8K and MATH, and include selected examples from MathInstruct \citep{MathInstruct}. Additionally, we incorporate the training portion of Lila-OOD \citep{lila}, in which problems have been addressed using CoT or PoT methodologies. This English dataset encompasses a wide array of mathematical disciplines including algebra, probability, number theory, calculus, and geometry.
    \item \textbf{Chinese mathematical datasets}: Our Chinese dataset consists of K-12 mathematical problems distributed across 76 different sub-topics, such as linear equations. Each problem is furnished with solutions detailed in both CoT and tool-integrated reasoning styles.
\end{itemize}

\subsection{Training and Evaluating DeepSeekMath-Instruct 7B}

This section presents DeepSeekMath-Instruct 7B, a model derived from DeepSeekMath-Base and further refined through instruction tuning specifically for mathematical tasks. During the training process, individual examples are randomly concatenated to fill up to a context window of 4K tokens. The model undergoes 500 training steps utilizing a batch size of 256 and a fixed learning rate set to 5e-5.

The mathematical capabilities of the models are assessed in settings both with and without tool assistance, across four benchmarks designed for quantitative reasoning in both English and Chinese. For comparison, our evaluation includes leading contemporaneous models:

\begin{itemize}[topsep=0pt]
    \item \textbf{Closed-source models} encompass: (1) the GPT series, with GPT-4 \citep{gpt4} and GPT-4 Code Interpreter~\footnote{\url{https://openai.com/blog/chatgpt-plugins\#\#code-interpreter}} standing out as the most advanced, (2) Gemini Ultra and Gemini Pro \citep{gemini}, (3) Inflection-2 \citep{inflection-2}, (4) Grok-1~\footnote{\url{https://x.ai/model-card}}, alongside several models launched by Chinese enterprises, including (5) Baichuan-3~\footnote{\url{https://www.baichuan-ai.com}}, and (6) the most recent version of GLM-4~\footnote{\url{https://open.bigmodel.cn/dev/api\#glm-4}}.

    from the GLM family \citep{glm}.

These models serve broad applications and have generally been refined through multiple alignment strategies.

\item \textbf{Open-source models} in this context comprise both general-purpose and math-specialized variants: (1) DeepSeek-LLM-Chat 67B \citep{deepseek-llm}, (2) Qwen 72B \citep{qwen}, (3) SeaLLM-v2 7B \citep{seallm}, and (4) ChatGLM3 6B \citep{chatglm3}, represent generic models, while models demonstrating advancements in mathematics include: (5) InternLM2-Math 20B~\footnote{\url{https://github.com/InternLM/InternLM-Math}}, an extension of InternLM2 augmented via mathematics pretraining followed by additional instruction tuning; (6) Math-Shepherd-Mistral 7B, which employs PPO-based fine-tuning \citep{schulman2017proximal} on Mistral 7B \citep{mistral} with supervision guided by a reward model evaluating intermediate reasoning steps; (7) the WizardMath suite \citep{wizardmath}, which enhances mathematical reasoning capabilities for Mistral 7B and Llama-2 70B \citep{llama2} through evolve-instruct (a methodology leveraging AI-generated instructions for instruction tuning) and PPO training on datasets predominantly sourced from GSM8K and MATH; (8) MetaMath 70B \citep{metamath}, which involves finetuning Llama-2 70B on expanded versions of GSM8K and MATH; (9) ToRA 34B \cite{tora}, created by adapting CodeLlama 34B to perform tool-assisted mathematical reasoning through targeted finetuning; and (10) MAmmoTH 70B \citep{MathInstruct}, produced by instruction-tuning Llama-2 70B using the MathInstruct collection.

As indicated in Table \ref{tab:sft_rl_math}, when evaluated without access to external tools, DeepSeekMath-Instruct 7B exhibits impressive capabilities in step-by-step reasoning tasks. In particular, on the challenging MATH competition dataset, our model achieves results that exceed those of all open-source alternatives and most proprietary systems (such as Inflection-2 and Gemini Pro) by a margin of at least 9\% absolute. This superiority holds even in comparison to models that are much larger in scale (for example, Qwen 72B) or models further optimized with specialized math-oriented reinforcement learning techniques (like WizardMath-v1.1 7B). Although DeepSeekMath-Instruct's performance is comparable to leading Chinese proprietary systems such as GLM-4 and Baichuan-3 on the MATH benchmark, it remains behind GPT-4 and Gemini Ultra.

In evaluation scenarios permitting the combination of natural language reasoning and program-driven tool application, DeepSeekMath-Instruct 7B achieves nearly 60\% accuracy on the MATH benchmark, outperforming every other available open-source model. Across additional benchmarks, our model demonstrates performance on par with DeepSeek-LLM-Chat 67B, the former leading model that is ten times its size.
\section{Reinforcement Learning}

\subsection{Group Relative Policy Optimization} Reinforcement learning (RL) has demonstrated significant potential in enhancing the mathematical reasoning capabilities of LLMs beyond what is achieved during the Supervised Fine-Tuning (SFT) phase \citep{wang2023math,wizardmath}. Here, we present Group Relative Policy Optimization (GRPO), our RL algorithm designed for both efficiency and effectiveness.

\subsubsection{From PPO to GRPO} Proximal Policy Optimization (PPO) \citep{schulman2017proximal} is a popular actor-critic reinforcement learning approach commonly adopted during the RL fine-tuning of large language models \citep{ouyang2022training}. This method enhances LLMs by seeking to maximize the following surrogate objective function:

\begin{equation}
\footnotesize
    \mathcal{J}_{PPO}(\theta) = \mathbb{E}{[q \sim P(Q), o \sim \pi_{\theta_{old}}(O|q)]} \frac{1}{|o|} \sum_{t=1}^{|o|} \min \left[ \frac{\pi_\theta(o_{t} | q, o_{<t})}{\pi_{\theta_{old}}(o_{t} | q, o_{<t})} A_{t}, \text{clip} \left( \frac{\pi_\theta(o_{t} | q, o_{<t})}{\pi_{\theta_{old}}(o_{t} | q, o_{<t})}, 1 - \varepsilon, 1 + \varepsilon \right)  A_{t} \right] ,
\end{equation}

Here, $\pi_{\theta}$ refers to the updated policy model, while $\pi_{\theta_{old}}$ denotes the previous version of the policy. The symbols $q$ and $o$ represent questions and the corresponding outputs, which are drawn from the dataset and the prior policy $\pi_{\theta_{old}}$. The parameter $\varepsilon$ serves as a clipping threshold in PPO, facilitating the stabilization of training dynamics. The term $A_t$ denotes the advantage, which is determined using Generalized Advantage Estimation (GAE) \citep{gae}, leveraging both the reward sequence $\{r_{\ge t}\}$ and a value function $V_{\psi}$ that is learned throughout training. As a result, PPO necessitates simultaneous training of a value function in addition to the policy model. To prevent excessive optimization of the reward model, it is common practice to introduce a per-token KL divergence penalty against a reference model into the reward computation at each token, following \citep{ouyang2022training}, namely,

\begin{equation}
    r_{t} = r_\varphi(q, o_{\le t}) - \beta \log\frac{\pi_{\theta}(o_{t}|q, o_{<t})}{\pi_{ref}(o_{t}|q, o_{<t})},
\label{eq:PPO-reward}
\end{equation}

where $r_\varphi$  is the reward model, $\pi_{ref}$ is the reference model, which is usually the initial SFT model, and $\beta$  is the coefficient of the KL penalty.

In PPO, the value function is generally implemented as an independent neural network, often comparable in scale to the policy network itself. This setup leads to significant computational and memory demands. Furthermore, the value function serves as a baseline for advantage estimation during RL training to mitigate variance. However, applying this framework to LLMs is challenging because, in this scenario, the reward model typically evaluates only the final token, complicating accurate value prediction at every token position. To overcome these challenges, we introduce Group Relative Policy Optimization (GRPO), as illustrated in Figure \ref{fig:grpo}. GRPO eliminates the requirement for an explicit value function, unlike PPO. Instead, it employs the mean reward across multiple generated responses to the same prompt as the baseline. Specifically, for a given prompt $q$, GRPO generates a set of responses $\{o_1, o_2, \cdots, o_G\}$ using the prior policy $\pi_{\theta_{old}}$, and updates the policy by maximizing the objective given below:

\begin{equation}
\footnotesize
\begin{split}
    \mathcal{J}_{GRPO}(\theta) &= \mathbb{E}{[q \sim P(Q), \{o_i\}_{i=1}^G \sim \pi_{\theta_{old}}(O|q)]}  \\
    & \frac{1}{G}\sum_{i=1}^G\frac{1}{|o_i|} \sum_{t=1}^{|o_i|} \left\{ \min \left[ \frac{\pi_\theta(o_{i,t} | q, o_{i,<t})}{\pi_{\theta_{old}}(o_{i,t} | q, o_{i,<t})} \hat{A}_{i,t}, \text{clip} \left( \frac{\pi_\theta(o_{i,t} | q, o_{i,<t})}{\pi_{\theta_{old}}(o_{i,t} | q, o_{i,<t})}, 1 - \varepsilon, 1 + \varepsilon \right)  \hat{A}_{i,t} \right] - \beta \mathbb{D}_{KL}\left[\pi_{\theta} || \pi_{ref}\right]\right\} ,
\end{split}
\label{eq:GRPO-obj}
\end{equation}

Here, $\varepsilon$ and $\beta$ serve as hyper-parameters, while $\hat{A}_{i,t}$ denotes the advantage, which is computed exclusively based on the relative rewards among outputs within each respective group; further explanation of this calculation is provided in the subsequent subsections. The approach of estimating advantages in a group-relative fashion within GRPO coheres naturally with the comparative framework characteristic of reward models, since such models are frequently optimized on comparison-based datasets generated from multiple answers to the same prompt. Additionally, it is important to observe that GRPO does not incorporate a KL penalty directly into the reward formulation. Instead, it introduces regularization by adding the KL divergence between the policy under training and the reference policy as a term in the loss function, thereby streamlining the computation of $\hat{A}_{i,t}$. Unlike the KL penalty utilized in (\ref{eq:PPO-reward}), our method estimates the KL divergence by employing the following unbiased estimator \citep{kl_approx}:

\begin{equation}
\small
    \mathbb{D}_{KL}\left[\pi_{\theta} || \pi_{ref}\right] = \frac{\pi_{ref}(o_{i,t}|q,o_{i,<t})}{\pi_{\theta}(o_{i,t}|q,o_{i,<t})}- \log\frac{\pi_{ref}(o_{i,t}|q,o_{i,<t})}{\pi_{\theta}(o_{i,t}|q,o_{i,<t})} - 1,
\end{equation}

which is guaranteed to be positive.

\begin{algorithm}[t]
  \small
  \caption{Iterative Group Relative Policy Optimization}
  \textbf{Input} Starting policy model $\pi_{\theta_{\text{init}}}$; reward models $r_\varphi$; collection of task prompts $\mathcal{D}$; 
  hyperparameters $\varepsilon$, $\beta$, $\mu$
  \begin{algorithmic}[1]
    \State Initialize policy model $\pi_\theta \leftarrow \pi_{\theta_{\text{init}}}$
    \For{iteration = 1, \dots, I}
       \State Set reference model $\pi_{ref} \leftarrow \pi_{\theta}$
      \For{step = 1, \dots, M}
      \State Draw a batch $\mathcal{D}_b$ from $\mathcal{D}$
      \State Assign $\pi_{\theta_{old}} \leftarrow \pi_{\theta}$
      \State For each prompt $q$ in $\mathcal{D}_b$, sample $G$ responses $\{o_i\}_{i=1}^G \sim \pi_{\theta_{old}} (\cdot \mid q) $
      \State Evaluate each sampled output $o_i$ with $r_{\varphi}$ to obtain rewards $\{r_i\}_{i=1}^{G}$
      \State For every $o_i$, estimate group-relative advantage $\hat{A}_{i,t}$ at token position $t$
      \For{GRPO iteration = 1, \dots, $\mu$}
        \State Optimize policy model $\pi_{\theta}$ via maximization of the GRPO objective (Equation \ref{eq:GC-GRPO})
      \EndFor
    \EndFor \State Update reward models $r_\varphi$ via continual learning, employing a replay buffer.
    \EndFor 
  \end{algorithmic}
  \textbf{Output} $\pi_\theta$
  \label{alg:iter-grpo}
\end{algorithm}

\subsubsection{Outcome Supervision RL with GRPO} Specifically, given a question $q$, the procedure involves sampling a set of outputs $\{o_1, o_2, \cdots, o_G\}$ from the previous policy $\pi_{\theta_{old}}$. Each of these outputs is evaluated by a reward model, which assigns a reward to each output, resulting in $G$ scores $\mathbf{r}=\{r_1, r_2, \cdots, r_G\}$. To ensure comparability, the rewards undergo normalization by subtracting the mean and dividing by the standard deviation within the group. Under the outcome supervision approach, the resulting normalized reward for each output $o_i$ is supplied only at the output's end, and all tokens associated with $o_i$ are assigned an advantage $\hat{A}_{i, t}$ equal to the normalized reward; mathematically, $\hat{A}_{i, t} = \widetilde{r}_i = \frac{r_i- {\rm mean}(\mathbf{r})}{{\rm std}(\mathbf{r})}$. The policy is then updated to maximize the target defined by equation (\ref{eq:GRPO-obj}).

\subsubsection{Process Supervision RL with GRPO}  
In outcome supervision, feedback is given only at the conclusion of each generated output, which is often inadequate and inefficient when guiding policies for intricate mathematical tasks. Inspired by \cite{wang2023math}, we adopt process supervision that assigns rewards at every step within a reasoning trajectory. Specifically, for a question $q$ and a set of $G$ sampled outputs $\{o_1, o_2, \cdots, o_G\}$, a process reward model evaluates each intermediate step, producing stepwise rewards: $\mathbf{R} = \{ \{r_1^{index(1)},\cdots,r_1^{index(K_1)}\}, \cdots,  \{r_G^{index(1)},\cdots,r_G^{index(K_G)}\} \}$, where $index(j)$ marks the position of the end token for the $j$-th step, and $K_i$ denotes the total step count for the $i$-th output. To stabilize the reward signal, we apply normalization using the overall mean and standard deviation: $\widetilde{r}_i^{index(j)} = \frac{r_i^{index(j)} - {\rm mean(\mathbf{R})}}{{\rm std(\mathbf{R})}}$.

Following this, for every token, the advantage is calculated by summing the normalized rewards of all subsequent steps, such that $\hat{A}_{i, t} = \sum_{index(j) \ge t} \widetilde{r}_i^{index(j)}$. The policy is then updated by maximizing the objective as given in equation (\ref{eq:GRPO-obj}).

\subsubsection{Iterative RL with GRPO} As the policy model advances during reinforcement learning, earlier versions of the reward model may no longer provide adequate supervision. To address this, we investigate an iterative approach to RL using GRPO. Algorithm \ref{alg:iter-grpo} outlines the process: in iterative GRPO, new datasets for the reward model are generated from samples produced by the current policy model. The existing reward model is then updated in a continual fashion through a replay strategy that retains 10\% of previous data. Subsequently, the policy model is assigned as the reference, and further training occurs using the updated reward model to guide the policy's learning.

\subsection{Training and Evaluating DeepSeekMath-RL}

Our reinforcement learning experiments utilize DeepSeekMath-Instruct 7B as the base model. The RL training dataset comprises chain-of-thought-style problems from GSM8K and MATH, selected from the SFT dataset and totaling approximately 144K samples. To specifically assess the influence of RL on benchmarks without exposure during training, other SFT-derived questions are omitted during RL. The methodology for constructing the reward model’s training data follows the approach of \citep{wang2023math}. The initial reward model is built on DeepSeekMath-Base 7B, with a learning rate set to $2\times10^{-5}$. In the context of GRPO, we assign the policy model a learning rate of $1\times10^{-6}$ and utilize a KL penalty coefficient of $0.04$. For every training question, we generate $64$ candidate completions. Sequence length is capped at $1024$ tokens, with a batch size of $1024$ per training step. The policy model receives a single parameter update after each exploration cycle. Evaluation of DeepSeekMath-RL 7B follows the same benchmarks as DeepSeekMath-Instruct 7B. For DeepSeekMath-RL 7B, tasks involving GSM8K and MATH with chain-of-thought prompting are categorized as in-domain, while all remaining evaluation benchmarks are considered out-of-domain.

Table \ref{tab:sft_rl_math} presents results for both open- and closed-source models, evaluated on English and Chinese benchmarks, considering both chain-of-thought and tool-assisted reasoning methods. Key findings include: 1) DeepSeekMath-RL 7B achieves accuracy rates of 88.2\% on GSM8K and 51.7\% on MATH using chain-of-thought strategies. These outcomes exceed those of all other open-source models within the 7B to 70B parameter range, and outperform most closed-source alternatives as well. 2) Notably, DeepSeekMath-RL 7B's training was exclusively limited to instruction-tuning data in a chain-of-thought format for GSM8K and MATH, initializing from DeepSeekMath-Instruct 7B. Despite this narrow training regime, DeepSeekMath-RL 7B consistently surpasses DeepSeekMath-Instruct 7B in all evaluated dimensions, underscoring the utility of reinforcement learning in this context.

\section{Discussion}
In this section, we will share our findings in pre-training and RL experiments. 

\subsection{Lessons Learnt in Pre-Training}

We begin by describing our observations from the pre-training phase. Except where noted, all training configurations follow those presented in Section \ref{sec:quality-policy}. Additionally, throughout this section, references to the DeepSeekMath Corpus pertain to a version containing 89B tokens, collected during the second round of data gathering.

\subsubsection{Code Training Benefits Mathematical Reasoning}

An oft-cited but unproven hypothesis posits that training with code enhances reasoning capabilities. In this work, we address this claim in the context of mathematics, providing evidence that code training bolsters models’ mathematical reasoning performance, irrespective of whether external tools are utilized.

To study how code training affects mathematical reasoning, we experimented with the following two-stage training and one-stage training settings:

\textbf{Two-Stage Training}

\begin{itemize}[topsep=0pt]
    \item \textbf{Code Training for 400B Tokens $\rightarrow$ Math Training for 150B Tokens}: DeepSeek-LLM 1.3B undergoes pretraining with 400B code tokens, after which it is further trained on 150B math tokens.
    \item \textbf{General Training for 400B Tokens $\rightarrow$ Math Training for 150B Tokens}: For comparison, we conduct an additional experiment where the model is pretrained on 400B general tokens (sampled from a comprehensive general domain dataset produced by DeepSeek-AI) rather than code tokens during the initial stage, with the goal of examining whether code tokens confer any benefits over general tokens in fostering mathematical reasoning skills.
\end{itemize}

\textbf{One-Stage Training}

\begin{itemize}[topsep=0pt]
    \item \textbf{Math Pretraining with 150B Tokens}: DeepSeek-LLM 1.3B undergoes pretraining on a dataset comprising 150B math-specific tokens;
    \item \textbf{Joint Training with 400B Code Tokens and 150B Math Tokens}: Sequentially pretraining on math after code data diminishes coding ability. Here, we explore if simultaneously introducing math and code tokens during a unified training stage can not only enhance the model’s proficiency in mathematical reasoning but also help prevent catastrophic forgetting.
\end{itemize}

\paragraph{Results}
Table \ref{tab:code-math} and Table \ref{tab:code-math-general-eval} demonstrate the downstream performance under different training settings.

Integrating code training into the curriculum improves program-aided mathematical reasoning, regardless of whether a two-stage or one-stage training framework is employed. According to Table \ref{tab:code-math}, solely incorporating code training during the initial stage of a two-stage protocol already substantially boosts performance on GSM8K and MATH datasets using Python as the solution tool. Introducing math training in the subsequent phase produces additional performance gains. Notably, within the one-stage training paradigm, blending code tokens with math tokens not only alleviates the problem of catastrophic forgetting associated with the two-stage approach but also enhances both coding capabilities (as demonstrated in Table \ref{tab:code-math-general-eval}) and program-aided mathematical reasoning (see Table \ref{tab:code-math}).

Training on code data also enhances mathematical reasoning abilities even when no tools are used. In the two-stage training paradigm, conducting code training in the first phase yields noticeable improvements. Furthermore, it accelerates and amplifies the gains during subsequent mathematical training, ultimately achieving superior results. In contrast, merging code and math tokens in a single-stage training process diminishes mathematical reasoning in the absence of tool usage. A possible explanation is that DeepSeek-LLM 1.3B, given its smaller model size, may not possess sufficient capacity to effectively integrate both code and mathematical datasets at the same time.

\subsubsection{ArXiv Papers Seem Ineffective in Improving Mathematical Reasoning}

ArXiv papers are frequently utilized as part of mathematical pre-training datasets \citep{minerva,gpt-f,llemma,mathpile}, yet systematic investigations into their influence on mathematical reasoning abilities remain scarce. Interestingly, our empirical findings reveal that incorporating arXiv material offers limited benefits for enhancing mathematical reasoning skills. We conduct experiments across models of various scales, specifically DeepSeek-LLM 1.3B and DeepSeek-Coder-Base-v1.5 7B \citep{deepseek-coder}, employing arXiv datasets processed through multiple distinct data preparation methods:

\begin{itemize}[topsep=0pt]
    \item \textbf{MathPile} \citep{mathpile}: a dataset comprising 8.9 billion tokens, curated using a set of heuristic cleaning and filtering strategies, with scientific articles from arXiv constituting more than 85\% of its content;
    \item \textbf{ArXiv-RedPajama} \citep{redpajama}: this corpus consists of all arXiv LaTeX documents after stripping out preambles, comments, macros, and reference sections, resulting in a dataset of 28.0 billion tokens.
\end{itemize}

For our experiments, we independently trained DeepSeek-LLM 1.3B on 150B tokens and DeepSeek-Coder-Base-v1.5 7B on 40B tokens, using arXiv-specific corpora for each model. The results suggest that arXiv papers contribute minimally to enhancing mathematical reasoning capabilities. Models trained exclusively on arXiv data did not exhibit measurable gains and, in some cases, showed declines across a spectrum of mathematical evaluation tasks utilized in this research. These tasks span quantitative reasoning datasets such as GSM8K and MATH (see Table \ref{tab:arxiv-cot}), multiple-choice evaluations including MMLU-STEM (Table \ref{tab:arxiv-cot}), and formal mathematics benchmarks like miniF2F (Table \ref{tab:arxiv_atp}).

However, this conclusion has its limitations and should be taken with a grain of salt. We have not yet studied:

\begin{itemize}[topsep=0pt]
    \item How arXiv tokens influence other math-centric tasks excluded from the current analysis, for instance, translating formal statements or proofs into their corresponding informal forms (informalization of theorems);
    \item The influence exerted by arXiv tokens when integrated with additional data sources;
    \item If the advantages conferred by including arXiv papers persist or become more evident as the scale of the model increases.
\end{itemize}

Thus, further exploration is required, which we leave for future studies.

\subsection{Insights of Reinforcement Learning}
\subsubsection{Moving Toward a Unified Paradigm}
In this section, we introduce a unified framework to systematically analyze various training approaches—including SFT, RFT, DPO, PPO, and GRPO—and proceed to empirically investigate the contributing factors within this paradigm. In general, the gradient of a training objective with respect to the parameters $\theta$ can be represented as:

\begin{equation}
    \nabla_{\theta}\mathcal{J}_{\textcolor{red}{\mathcal{A}}}(\theta) = \mathbb{E}[\underbrace{(q,o) \sim \textcolor{red}{\mathcal{D}}}_{Data \ Source}]\left( \frac{1}{|o|} \sum_{t=1}^{|o|}  \underbrace{GC_{{\mathcal{A}}}(q, o, t, \textcolor{red}{\pi_{{rf}}})}_{Gradient \ Coefficient}  \nabla_{\theta}\log \pi_{\theta}(o_t | q, o_{<t})\right).
\label{eq:objective}
\end{equation}

Three primary components are involved: 1) \textit{Data Source $\mathcal{D}$}, specifying the dataset used during training; 2) \textit{Reward Function $\pi_{{rf}}$}, providing the training reward signal; and 3) \textit{Algorithm $\mathcal{A}$}, which utilizes both the training data and reward signal to compute the gradient coefficient $GC$, thereby dictating the extent of penalization or reinforcement applied to the data. We examine a range of prototypical methods within this unified framework:

\begin{itemize}
    \item  \textbf{Supervised Fine-tuning (SFT)}: SFT involves adapting a pretrained model to specific data curated by humans through supervised learning.
    \item \textbf{Rejection Sampling Fine-tuning (RFT)}: RFT further adapts the SFT-trained model by fine-tuning it on selected outputs, which are generated by the SFT model for SFT prompts and filtered based on answer accuracy.
    \item \textbf{Direct Preference Optimization (DPO)}: DPO enhances the SFT model by training it on a set of expanded outputs, sampled from the SFT model, and using a pairwise DPO loss objective.
    \item \textbf{Online Rejection Sampling Fine-tuning (Online RFT)}: Unlike standard RFT, Online RFT begins with a policy model initialized from the SFT model, then continues to improve it by fine-tuning on augmented outputs sampled directly from the evolving policy during training.
    \item \textbf{PPO/GRPO}: PPO/GRPO also start with the SFT model to initialize the policy, and subsequently use reinforcement learning, sampling outputs from the real-time policy for further training.
\end{itemize}

We summarize the components of these methods in Table \ref{tab:data_policy}.
Please refer to Appendix \ref{app:analysis-rl} for a more detailed derivation process.

\paragraph{Observation about Data Source} The origin of the data is categorized into two primary types: online sampling and offline sampling. In the case of online sampling, the data used for training is collected from the ongoing explorations of the policy model as it is being trained. Conversely, offline sampling utilizes data generated from the initial SFT model prior to any interactive training. RFT and DPO rely on offline sampling, whereas Online RFT and GRPO employ data acquired through online sampling.

As illustrated in Figure \ref{fig:rl-analysis}, our results indicate that Online RFT achieves markedly better performance than RFT across both benchmarks. Notably, while Online RFT and RFT perform similarly during the initial training phase, Online RFT demonstrates clear superiority as training progresses. This outcome aligns with expectations: at the outset, the actor and the SFT model are highly similar, with sampled trajectories differing only slightly. Over time, however, the actor’s samples increasingly diverge from those of the SFT model, making real-time data collection substantially more beneficial during later training stages.

\paragraph{Observation about Gradient Coefficient} In this framework, the reward signal is converted into a gradient coefficient that directs model parameter updates. Our experimental setup distinguishes between two forms of reward function: `Rule' and `Model'. The `Rule' approach evaluates response quality purely by factual correctness, while the `Model' approach involves training a reward model that assigns scores to responses, utilizing data labeled according to rule-based judgments. As illustrated in Equations \ref{eq:GC-RFT} and \ref{eq:GC-GRPO}, a fundamental distinction exists between GRPO and Online RFT: GRPO adapts its gradient coefficient in accordance with the reward model's outputs, thereby enabling the model to reinforce or penalize responses in proportion to their respective reward magnitudes. On the other hand, Online RFT does not offer such nuanced feedback—it treats all correct responses identically and refrains from penalizing incorrect answers.

As illustrated in Figure \ref{fig:rl-analysis}, GRPO outperforms online RFT, emphasizing the advantages of adjusting the coefficients for positive and negative gradients. Moreover, GRPO+PS achieves better results than GRPO+OS, demonstrating the value of employing more granular, step-specific gradient coefficients. We also investigate the impact of iterative RL, implementing two iterations in our experimental setup. Figure \ref{fig:iter-rl} reveals that iterative RL produces notable performance enhancements, with the most significant gains occurring during the initial iteration.

\subsubsection{Why RL Works?} In this study, we apply reinforcement learning using a selected portion of instruction tuning data and observe a notable improvement in model performance over the instruction-tuned baseline. To elucidate the mechanism behind the efficacy of reinforcement learning, we compare Pass@K and Maj@K accuracies for both the Instruct and RL models across two evaluation benchmarks. According to the results presented in Figure \ref{fig:maj-pass}, reinforcement learning leads to a marked gain in Maj@K, whereas Pass@K remains largely unaffected. These observations suggest that RL strengthens the model's output by making the answer distribution more stable or consistent, specifically, \textbf{implying that performance gains arise from promoting correct responses within the TopK outputs rather than advancing the model’s core competencies.} Analogously, \citep{wang2023making} reports a \textbf{misalignment problem} encountered during reasoning tasks in SFT models, and demonstrates that the application of preference alignment techniques can enhance the reasoning abilities of SFT models \citep{yuan2023rrhf,song2023preference,wang2023making}.

\subsubsection{How to Achieve More Effective RL?}
\label{sec:effec_rl}
Our results indicate that RL is highly effective when applied to mathematical reasoning problems. Additionally, we introduce a cohesive framework that brings various significant training approaches under a single perspective. Within this unified view, all strategies can be interpreted as variations of either standard or reduced RL approaches. As outlined in Equation \ref{eq:objective}, three fundamental elements are identified: Data Source, Algorithm, and Reward Function. Below, we discuss potential avenues for advancing these three aspects in future research.

\paragraph{Data Source} The data source constitutes the foundational input for all training paradigms. Within the scope of RL, we define the data source as unlabeled questions accompanied by responses generated by the policy model. For the present study, the dataset is limited to questions originating from the instruction tuning phase, with outputs produced via simple nucleus sampling. We hypothesize that this design choice may be partially responsible for the RL pipeline's effectiveness being confined to improvements in Maj@K. Moving forward, we intend to evaluate our RL pipeline using out-of-distribution question prompts and to incorporate \textbf{advanced sampling (decoding) strategies}, such as those informed by tree-search algorithms \citep{yao2023tree}. In addition, the use of \textbf{efficient inference techniques} \citep{xia-etal-2023-speculative,leviathan2023fast,kwon2023efficient,xia2024unlocking}, which critically influence how effectively policy models can explore, will also be a focal point due to their considerable impact.

\paragraph{Algorithms} Algorithms utilize both the data and reward signals to compute the gradient coefficient, which in turn is used to adjust the model parameters. Drawing from Equation \ref{eq:objective}, contemporary approaches generally place full \textbf{TRUST} in the reward function's feedback when modulating the conditional probability of a specific token. Nevertheless, there is no guarantee that the reward signals maintain their reliability at all times, particularly when dealing with highly complex tasks. As an illustration, the PRM800K datasets \citep{lightman2023let}, despite their extensive annotation process conducted by skilled annotators, still exhibit around 20\% of incorrectly labeled instances\footnote{\url{https://github.com/openai/prm800k/issues/12\#issuecomment-1728491852}}. Therefore, it is crucial to investigate reinforcement learning algorithms capable of withstanding noisy reward signals. In this context, we anticipate that alignment strategies based on the \textbf{WEAK-TO-STRONG} paradigm \citep{burns2023weak} hold the potential to profoundly reshape learning algorithm design.

\paragraph{Reward Function} The reward function provides the main training signal in RL, and is typically implemented as a neural reward model. We identify three key avenues for advancing reward models: 1) \textbf{Improving the reward model’s generalization capability.} It is critical that the reward model be able to generalize well to both out-of-distribution queries and sophisticated generation outputs; failure to do so could result in reinforcement learning merely reinforcing the existing output distribution of LLMs rather than enhancing their core competencies; 2) \textbf{Capturing the uncertainty inherent in the reward model.} Accounting for uncertainty may facilitate connections between suboptimal reward models and learning strategies designed to improve them; 3) \textbf{Developing process reward models with high efficiency and quality} that yield fine-grained signals to guide the reasoning process during training \citep{lightman2023let,wang2023math}.

\section{Conclusion, Limitation, and Future Work} This work introduces DeepSeekMath, a model that achieves superior results compared to all open-source counterparts on the competition-level MATH benchmark and demonstrates performance that rivals proprietary systems. DeepSeekMath~was built by initializing from DeepSeek-Coder-v1.5 7B and subsequently trained continuously on 500B tokens, incorporating 120B math-specific tokens curated from Common Crawl as a substantial part of its dataset. Through a detailed ablation analysis, we reveal that content from web pages constitutes a promising resource for obtaining high-quality mathematical material, whereas the utility of arXiv data fell short of our initial assumptions. We propose Group Relative Policy Optimization (GRPO)—an adaptation of Proximal Policy Optimization (PPO)—which enhances mathematical reasoning proficiency while reducing memory requirements. Our experiments confirm that GRPO remains beneficial even when DeepSeekMath-Instruct 7B performs strongly across benchmark tests. Moreover, we offer a unified framework for interpreting various approaches and identify several promising avenues for advancing reinforcement learning methods.

While DeepSeekMath~demonstrates strong performance on quantitative reasoning tasks, its abilities in geometry and theorem-proving are notably less robust compared to proprietary models. For example, in our preliminary evaluation, the model was unable to address questions involving triangles and ellipses, which suggests that there may be a bias in the data used during pre-training and fine-tuning stages. Moreover, due to the model’s current scale, DeepSeekMath~lags behind GPT-4 when it comes to few-shot learning. While GPT-4 shows improvement when given few-shot examples, DeepSeekMath~performs nearly identically under zero-shot and few-shot conditions. Moving forward, we plan to enhance our curated data selection process to build a higher-quality pre-training dataset. Furthermore, we intend to investigate additional methods for optimizing reinforcement learning in LLMs, as discussed in Section \ref{sec:effec_rl}.

\end{document}